% SPDX-License-Identifier: CC-BY-4.0
\section{Related work}

\paragraph{Systems and standards.}
The setting is no longer notional. ESA's Moonlight/LCNS has a signed prime
contract and a space segment in build \citep{esa_moonlight,telespazio2024moonlight}.
NASA, ESA and JAXA share the Augmented Forward Signal under version 5 of the
LunaNet Interoperability Specification \citep{nasa2025lnis,nasa2025afs}. ESA has
adopted NovaMoon, a south-pole reference station that combines VLBI, ranging,
retroreflectors and a GNSS receiver \citep{molli2026novamoon}. And GNSS has
actually been received on the surface: LuGRE took the first lunar-surface fix on
Blue Ghost Mission~1 \citep{lugre2026}, with SSTL's Lunar Pathfinder and its
NaviMoon receiver to follow (launch not earlier than 2026).

\paragraph{Frames, time and tracking.}
The IAU's 2024 General Assembly defined a lunar coordinate time scale (TCL) and a
lunar reference system \citep{iau2024lcrs}; the relativistic Earth--Moon rate is
fixed by a published framework \citep{ashby2024lunartime}, and a realisation has
been worked out in the Moonlight context
\citep{fienga2024moonlight,bourgoin2026lunartime}. Body-fixed orientation follows
the IAU/WGCCRE conventions \citep{archinal2018wgccre} over the JPL DE440
ephemeris \citep{park2021de440}. The Earth-baseline geodetic tie to a lunar
beacon is the descendant of deep-space $\Delta$-DOR, the one-nanoradian
measurement system of the DSN \citep{curkendall2013deltador}.

\paragraph{Estimation and tooling.}
Lunar orbit-and-clock estimation and constellation design are active areas
\citep{iiyama2026odclock,bhamidipati2023constellation}, and recent lower-bound
work finds that a surface reference station buys sub-metre accuracy even under
sparse satellite visibility \citep{pohlmann2025hybrid}, which lines up with the
differential and observability results we report below. There is no shortage of
open orbit-determination software either: GMAT \citep{nasa_gmat}, Orekit
\citep{orekit}, RTKLIB \citep{takasu2009}, the Rust toolkit Nyx/ANISE
\citep{nyxspace2025}, and, for the lunar case specifically, the LuPNT simulator
\citep{iiyama2023lupnt}; kshana is the broader engine these results sit in
\citep{baweja2026kshana}. What we add is not a new estimator. It is an open,
reproducible observability finding for the lunar joint-OD problem that anyone can
regenerate from a committed scenario, with every figure labelled validated or
modelled by the engine itself, and with the information-not-count control that
keeps a geometric claim from being mistaken for a data-volume effect. We have not
found that combination in either the open lunar simulators or the closed
mission-analysis studies.
